% part: intuitionistic-logic
% chapter: introduction
% section: constructive-reasoning

\documentclass[../../../include/open-logic-chapter]{subfiles}

\begin{document}

\olfileid[zh]{int}{int}{cr}

\section{构造性推理}

与通过模态算子或二阶量词扩张经典逻辑不同，直觉主义逻辑之所以是「非经典的」（non-classical），在于它对经典逻辑作了限制。经典逻辑在许多方面是\emph{非构造性的}（non-constructive）。直觉主义逻辑旨在刻画一种更具「构造性」的推理，这种推理是某类构造性数学的特征。下面的例子可以用来说明其中一些动机。

设有人声称他们确定了一个自然数~$n$，它满足如下性质：若 $n$ 是偶数，则黎曼猜想为真；若 $n$ 是奇数，则黎曼猜想为假。这可是个天大的消息！{}黎曼猜想是否成立是数学中重大的开放问题之一，而他们似乎把这个问题化约为一次计算，也就是化约为判断某个特定的数是否为偶数。

那么~$n$ 究竟是多少呢？他们这样描述：$n$ 是等于 $2$ 的自然数（若黎曼猜想为真），否则等于 $3$。

你愤怒地要求退款。从经典的立场看，上述描述确实确定了一个唯一的 $n$ 值；但你所真正想要的是一个能够\emph{显明地}（explicitly）给出的 $n$ 值。

再举一个也许不那么牵强的例子，考虑如下问题。我们知道，可以把一个无理数提高到一个有理数次幂，而得到有理结果。例如，$\sqrt{2}^2 = 2$。不太清楚的是，能否把一个无理数提高到某个\emph{无理数}次幂，而得到有理结果。如下定理对此作出了肯定回答：

\begin{thm}
存在无理数 $a$ 和 $b$，使得 $a^b$ 是有理数。
\end{thm}

\begin{proof}
考虑 $\sqrt{2}^{\sqrt{2}}$。若它是有理的，则我们已完成证明：可令 $a = b = \sqrt{2}$。否则，它就是无理的。于是我们有
\[
(\sqrt{2}^{\sqrt{2}})^{\sqrt{2}} = \sqrt{2}^{\sqrt{2} \cdot
  \sqrt{2}} = \sqrt{2}^2 = 2,
\]
而这是有理的。所以在这种情况下，令 $a$ 为
$\sqrt{2}^{\sqrt{2}}$，令 $b$ 为~$\sqrt 2$ 即可。
\end{proof}

这构成一个有效的证明吗？大多数数学家觉得是的。但这里又有一点令人不太满意：我们证明了具有某种性质的一对实数存在，却无法说出它究竟是\emph{哪一对}数。我们也可以用另一种方式来证明同一结果，使得这一对 $a$、$b$ \emph{确实}在证明中被给出：取 $a = \sqrt{3}$、$b = \log_3 4$。于是
\[
a^b = \sqrt{3}^{\log_3 4} = 3^{1/2 \cdot \log_3 4} = (3^{\log_3
  4})^{1/2} = 4^{1/2}= 2,
\]
因为 $3^{\log_3 x} = x$。

直觉主义逻辑旨在刻画这样一种推理，在这种推理中像第一个证明里那样的步骤是不被允许的。要证明某个满足~$!A(x)$ 的 $x$ 存在，你就必须给出一个具体的~$x$，以及它满足 $!A$ 的证明，就像第二个证明那样。要证明 $!A$ 或 $!B$ 成立，你就必须能够证明其中某一个成立。

形式地说，直觉主义逻辑就是在某一方面限制经典逻辑的!!a{derivation}系统所得到的东西。从数学的观点看，它们只是一些形式演绎系统，但如前所述，它们旨在刻画某种数学推理。我们可以把它看作在某种数学哲学观（如 Brouwer 的直觉主义）下得以辩护的那种推理；也可以把它看作一种更具「具体性」且令人满意的数学推理（沿着 Bishop 的构造主义思路）；同时，我们也可以争论这一形式描述是否刻画了非形式的动机。但无论我们持何种哲学立场，我们都可以把直觉主义逻辑当作一种以形式方式呈现的逻辑来研究；而且无论出于何种理由，许多数理逻辑学家都觉得这样做很有意思。

\end{document}
